\section{Escopo}
\label{sec_Escopo}


Conforme ilustrado na Figura~\ref{fig_etapas_rvdb}, uma missão de \gls{rvdb} pode ser descrita como uma sequência de etapas que conduzem a espaçonave perseguidora desde a inserção orbital até a aproximação final em relação ao alvo. 
Essa decomposição por fases fornece uma referência operacional para a organização das etapas consideradas na simulação desenvolvida neste estudo.

Com base nessa estrutura, são definidas as fases analisadas ao longo do procedimento de simulação, incluindo as etapas de aproximação de longa distância, aproximação de curta distância e aproximação final. 
Para cada fase são apresentados os principais vínculos operacionais, as hipóteses adotadas e os modelos dinâmicos utilizados na representação do sistema.

Em particular, na etapa de injeção orbital, considera-se o requisito de compatibilidade geométrica entre as órbitas do perseguidor e do alvo, com ênfase na condição de coplanaridade. 
Essa condição influencia o planejamento da janela de lançamento e reduz a necessidade de manobras posteriores de correção de trajetória.
